\subsection{群胚的同伦等价态射}

在本节中，我们将用记号 \(u \sim v\) 来表示两个群胚态射 u 与 v 同伦。

定义 2. —— 设 G, G' 为群胚，\(\varphi\) 为从 G 到 G' 的态射。\(\varphi\) 的同伦逆是指从 G' 到 G 的群胚态射 \(\psi\) ，使得态射 \(\psi \circ \varphi\) 与 \(\varphi \circ \psi\) 分别同伦于 \(\mathrm{Id}_{\mathrm{G}}\) 与 \(\mathrm{Id}_{\mathrm{G}'}\) 。称 \(\varphi\) 为同伦等价态射，如果存在 \(\varphi\) 的同伦逆。

群胚的同构是同伦等价态射。

设 G 与 G' 为群胚。设 \(\varphi\) 与 \(\varphi'\) 为从 G 到 G' 的、互相同伦的群胚态射。若 \(\varphi\) 是同伦等价态射，则 \(\varphi'\) 亦然。事实上，若以 \(\psi\) 表示 \(\varphi\) 的一个同伦逆，则有 \(\psi \circ \varphi' \sim \psi \circ \varphi \sim \mathrm{Id}_{\mathrm{G}}\) 与 \(\varphi' \circ \psi \sim \varphi \circ \psi \sim \mathrm{Id}_{\mathrm{G}'}\) ，这就证明 \(\psi\) 是 \(\varphi'\) 的同伦逆。

设 G, G' 与 G'' 为群胚，\(\varphi\colon\mathrm{G}\to\mathrm{G}'\) 、\(\varphi'\colon\mathrm{G}'\to\mathrm{G}''\) 、\(\psi\colon\mathrm{G}'\to\mathrm{G}\) 、\(\psi'\colon\mathrm{G}''\to\mathrm{G}'\) 为群胚态射。那么，在下列条件中：

\begin{enumerate}
\def\labelenumi{(\roman{enumi})}
\item
  \(\psi\) 是 \(\varphi\) 的同伦逆；
\item
  \(\psi'\) 是 \(\varphi'\) 的同伦逆；
\item
  \(\psi \circ \psi'\) 是 \(\varphi' \circ \varphi\) 的同伦逆；
\end{enumerate}

任何两个蕴含第三个。事实上，先设 (i) 与 (ii) 成立；则有

\[
\psi \circ \psi^ {\prime} \circ \varphi^ {\prime} \circ \varphi \sim \psi \circ \mathrm{Id} _ {\mathrm{G} ^ {\prime}} \circ \varphi \sim \psi \circ \varphi \sim \mathrm{Id} _ {\mathrm{G}}
\]

且类似地有 \(\varphi' \circ \varphi \circ \psi \circ \psi' \sim \mathrm{Id}_{\mathrm{G}'}\) ，由此得 (iii)。

\[
\varphi^ {\prime} \circ \psi^ {\prime} \sim \varphi^ {\prime} \circ \varphi \circ \psi \circ \psi^ {\prime} \sim \mathrm{Id} _ {\mathrm{G} ^ {\prime \prime}}
\]

以及

\[
\psi^ {\prime} \circ \varphi^ {\prime} \sim (\varphi \circ \psi) \circ \psi^ {\prime} \circ \varphi^ {\prime} \circ (\varphi \circ \psi) \sim \varphi \circ \psi \sim \mathrm{Id} _ {\mathrm{G} ^ {\prime}},
\]

由此得条件 (ii)。条件 (ii) 与 (iii) 蕴含条件 (i) 的证明是类似的。

特别地，若态射 \(\varphi\) 、\(\varphi'\) 、\(\varphi'\circ\varphi\) 中有两个是同伦等价态射，则第三个亦然。

命题 1. —— 设 G, G' 为群胚，\(\varphi\) 为从 G 到 G' 的态射，A 为与 G 的每个轨道都相交的 G 的顶点集的子集。为使 \(\varphi\) 是同伦等价态射，必须且只需下列条件成立：

\begin{enumerate}
\def\labelenumi{(\roman{enumi})}
\item
  映射 \(\mathrm{Orb}(\varphi)\) （从 \(\mathrm{Orb}(\mathrm{G})\) 到 \(\mathrm{Orb}(\mathrm{G}')\) ，由 \(\varphi\) 通过取轨道而得）是双射；
\item
  对一切 \(a \in \mathrm{A}\) ，同态 \(\varphi_a \colon \mathrm{G}_a \to \mathrm{G}_{\varphi(a)}'\) 是双射的。
\end{enumerate}

先设 \(\varphi\) 是同伦等价态射，并设 \(\psi\) 为 \(\varphi\) 的一个同伦逆。则

\[
\operatorname{Orb} (\psi) \circ \operatorname{Orb} (\varphi) = \operatorname{Orb} (\psi \circ \varphi) = \operatorname{Orb} (\operatorname{Id} _ {\mathrm{G}}) = \operatorname{Id} _ {\operatorname{Orb} (\mathrm{G})},
\]

因为互相同伦的群胚态射通过取轨道诱导相同的映射。同样，\(\mathrm{Orb}(\varphi) \circ \mathrm{Orb}(\psi) = \mathrm{Id}_{\mathrm{Orb}(G')}\) 。因此映射 \(\mathrm{Orb}(\varphi)\) 是双射，由此得断言 (i)。

\[
\psi_ {\varphi (a)} \circ \varphi_ {a} = (\psi \circ \varphi) _ {a};
\]

由于 \(\psi \circ \varphi\) 同伦于 \(\mathrm{Id}_{\mathrm{G}}\) ，同态 \((\psi \circ \varphi)_a\) 是双射的（参看第 180 页），因此 \(\varphi_a\) 是单射的而 \(\psi_{\varphi(a)}\) 是满射的。交换 \(\varphi\) 与 \(\psi\) 的角色，即见 \(\varphi_a\) 也是满射的，由此得条件 (ii)。

现在设条件 (i) 与 (ii) 成立，我们来证明 \(\varphi\) 是同伦等价态射。

先处理 G' 的每个轨道退缩为一点的情形。

对每个顶点 \(b\) （\(G'\) 的），选取顶点 \(u(b)\) （\(G\) 的、属于 \(A\) 的），使其在 \(\varphi\) 下的像为 \(b\) 。这是可能的，因为映射 \(\operatorname{Orb}(\varphi)\) 是满射的，且 \(A\) 与 \(G\) 的每个轨道相交。设 \(f\) 为 \(G'\) 的箭；按假设有 \(o(f) = t(f)\) ；令 \(b = o(f)\) 且 \(a = u(b)\) 。由条件 (ii)，存在唯一的箭 \(v(f) \in G_a\) ，它在 \(\varphi\) 下的像为 \(f\) 。偶 \(\psi = (u, v)\) 是从 \(G'\) 到 \(G\) 的群胚态射。我们来证明 \(\psi\) 是 \(\varphi\) 的同伦逆。我们已有 \(\varphi \circ \psi = \mathrm{Id}_{\mathrm{G}'}\) （按 \(\psi\) 的构造）。

设 \(x\) 为 G 的顶点。令 \(a = \psi(\varphi(x))\) ；这是 A 中满足 \(\varphi(a) = \varphi(x)\) 的元素。由于 \(\mathrm{Orb}(\varphi)\) 是单射的，\(a\) 属于 \(x\) 在 G 中的轨道，且存在 G 中的箭 \(f\) 连接 \(a\) 与 \(x\) 。于是箭 \(h(x) = \psi(\varphi(f))^{-1}f\) 连接 \(a\) 与 \(x\) ，且有 \(\varphi(h(x)) = e_{\varphi(a)}\) 。

我们来证明这样定义的映射 \(h \colon \operatorname{Som}(\mathbf{G}) \to \operatorname{Fl}(\mathbf{G})\) 是连接 \(\psi \circ \varphi\) 与 \(\operatorname{Id}_{\mathbf{G}}\) 的同伦。按构造，定义 1 的条件 (i) 成立。设 \(f\) 为 \(\mathbf{G}\) 的箭，\(x\) 为其起点，\(y\) 为其靶。有 \(\varphi(x) = \varphi(y)\) ；令 \(a = \psi(\varphi(x)) = \psi(\varphi(y))\) 。箭 \(h(x)fh(y)^{-1}\) 与 \(\psi \circ \varphi(f)\) 都属于 \(\mathrm{G}_a\) ，且都有像 \(\varphi(f)\) （在 \(\mathrm{G}'_{\varphi(a)}\) 中）。由于映射 \(\varphi_a\) 是单射的，有 \(h(x)fh(y)^{-1} = \psi \circ \varphi(f)\) 。于是定义 1 的条件 (ii) 成立。

现在来证明一般情形下的命题 1。设 X 为 G' 的极大定向森林（II, 第 157 页, 命题 1）。设 G'' 为由 G' 经缩并 X 的箭而得的群胚，\(\varphi'\colon G' \to G''\) 为典范态射。态射 \(\varphi'\) 满足命题 1 的条件（II, 第 170 页, 注 1 及第 178 页, 推论 2），态射 \(\varphi'\circ\varphi\) 亦然。

由于 G'' 的轨道都退缩为点，由已证的特例推出 \(\varphi'\) 与 \(\varphi'\circ\varphi\) 都是同伦等价态射，因而 \(\varphi\) 也是。

推论 1. —— 设 G 为群胚，A 为集合，\(f\colon A \to \operatorname{Som}(G)\) 为映射。若 \(f\) 的像与 G 的每个轨道相交，则典范群胚态射 \(\varphi\) （从 \(f^{*}G\) 到 G）是同伦等价态射。

按逆像群胚的定义（II, 第 166 页, 例 4），A 是群胚 \(f^{*}\mathrm{G}\) 的顶点集，且有 \(\mathrm{Fl}_{a,b}(f^{*}\mathrm{G}) = \mathrm{Fl}_{f(a),f(b)}(\mathrm{G})\) （对 A 的每一对元素 \((a,b)\) ）。此外，有 \(\operatorname{Som}(\varphi) = f\) ，且 \(\operatorname{Fl}(\varphi)\) 诱导从 \(\operatorname{Fl}_{a,b}(f^{*}\mathrm{G})\) 到 \(\operatorname{Fl}_{f(a),f(b)}(\mathrm{G})\) 的恒同映射。因此映射 \(\operatorname{Orb}(\varphi)\) 是双射的，且同态 \(\varphi_{a}\colon (f^{*}\mathrm{G})_{a} \to \operatorname{G}_{f(a)}\) 是同构（对一切 \(a \in A\) ）。于是命题 1 的假设得到满足。

推论 2. —— 设 G 为群胚，X 为 G 的定向森林，G' 为由 G 经缩并 X 的箭而得的群胚。则从 G 到 G' 的典范态射是同伦等价态射。

这由 II, 第 170 页, 注 1 及 II, 第 178 页, 推论 2 得出，它们表明命题 1 的假设得到满足。

